% Rúbrica: Estrategia (15 puntos): accionables, descripción de usuarios y usabilidad del modelo.
% Fuente: las capturas de docs/capturas/documento/ (provisionales: se regeneran desde el tag de entrega),
% frontend/src/app/ y las cifras de data-v3 que da generar_figuras.py (cifras_del_recorrido).
\section{Estrategia}\label{sec:estrategia}

\subsection{Usuarios y necesidades}\label{sec:estrategia-usuarios}

Quien usa NexPlay está por comprar un juego en Steam desde México. Necesita tres cosas: saber
antes de pagar si el juego tiende a decepcionar pronto, entender por qué, y decidir entre opciones
con eso a la vista. El sitio le pide un perfil corto, que se declara en un formulario: compras al
año, horas por semana, tolerancia a la fricción, las etiquetas de Steam que prefiere o rechaza y la
plataforma. Ese perfil sirve para la afinidad, que dice qué tanto encaja un juego con quien lo
declaró; no mueve el riesgo, que es del título (sección~\ref{sec:modelacion-bmas}).

\subsection{Accionables}\label{sec:estrategia-accionables}

Cada resultado del proyecto lleva a una acción concreta en el sitio:
\begin{itemize}
  \item \textbf{Comparar antes de decidir.} Comparar pone lado a lado la banda de riesgo, el motivo
  principal, el precio y la crítica de los juegos elegidos (figura~\ref{fig:accionables}a).
  \item \textbf{Revisar si las quejas son de las que le afectan.} Los motivos dicen de qué se
  quejan las reseñas con la señal; un juego cuyas quejas son de rendimiento pesa distinto para
  quien tiene un equipo de sobra que para quien no (figura~\ref{fig:recorrido}e y f).
  \item \textbf{Preguntarle a Nia,} que contesta con los datos del catálogo y explica la banda de
  un juego (sección~\ref{sec:interfaz}).
  \item \textbf{Planear las dos primeras horas.} El Inicio recuerda la ventana de reembolso y lo
  que más mencionan las negativas tempranas, para probar eso primero mientras el juego todavía se
  puede devolver; Steam decide cada solicitud, y el sitio lo dice (figura~\ref{fig:accionables}c).
  \item \textbf{Buscar alternativas del mismo género con riesgo bajo.} Explorar agrupa el catálogo
  en estantes por banda y se filtra por género (figura~\ref{fig:accionables}b).
\end{itemize}

\subsection{Recorrido: PAYDAY 3 contra Dead Space}\label{sec:estrategia-recorrido}

Una persona duda entre PAYDAY 3 y Dead Space, dos juegos de acción con precios casi iguales al corte
de datos: \cifraPaydayPrecio{} pesos PAYDAY 3, al \cifraPaydayFechaPrecio{}, y
\cifraDeadSpacePrecio{} Dead Space, al \cifraDeadSpaceFechaPrecio{}, con \cifraDeadSpaceDescuento{}
de descuento sobre \cifraDeadSpacePrecioInicial{}. Son los precios que habría visto quien comprara
esos días. El recorrido tiene cinco pasos:
\begin{enumerate}
  \item En el Inicio busca los dos. La ficha de PAYDAY 3 dice riesgo alto y la de Dead Space, riesgo
  bajo, cada una con su frase (figura~\ref{fig:recorrido}a y b). Dead Space es uno de los
  \cifraJuegosExternos{} juegos que el modelo no vio al entrenar.
  \item «Qué mueve esta estimación» da los factores (figura~\ref{fig:recorrido}c y d). PAYDAY 3
  tiene una nota de Metacritic de \cifraPaydayMetacritic{}, por debajo del promedio del catálogo, y
  eso sube su riesgo estimado; Dead Space tiene \cifraDeadSpaceMetacritic{}. El descuento de Dead
  Space aparece como un factor que lo baja, pero con evidencia débil, así que no es parte del
  argumento.
  \item Los motivos marcan la diferencia (figura~\ref{fig:recorrido}e y f). En PAYDAY 3, la queja
  principal es el contenido, en el \cifraPaydayMotivoPrincipal{} de las
  \cifraPaydayClasificadas{} reseñas clasificadas de sus \cifraPaydayCasos{} con la señal. En Dead
  Space es el rendimiento, en el \cifraDeadSpaceMotivoPrincipal{} de \cifraDeadSpaceClasificadas{},
  de \cifraDeadSpaceCasos{}. Las quejas de rendimiento pesan según el equipo de quien juega; las de
  contenido, no.
  \item Comparar los pone lado a lado (figura~\ref{fig:accionables}a), y Nia puede explicar la
  banda de cualquiera de los dos.
  \item Si quiere otro juego de acción con riesgo bajo, Explorar muestra un estante de
  \cifraEstanteAccionBajo{} (figura~\ref{fig:accionables}b). Va del score más bajo al más alto, y
  el precio 0 empuja al frente a los gratis y a los que no tienen precio: Grand Theft Auto V Legacy
  lleva el aviso de precio imputado, y Team Fortress 2 y Apex Legends, el de juego gratis. La
  alternativa sin aviso es Resident Evil 4.
\end{enumerate}

\begin{figure}[htbp]
  \centering
  \begin{subfigure}{0.49\textwidth}\includegraphics[width=\linewidth]{capturas/02-payday-veredicto.png}\caption{}\end{subfigure}\hfill
  \begin{subfigure}{0.49\textwidth}\includegraphics[width=\linewidth]{capturas/04-deadspace-veredicto.png}\caption{}\end{subfigure}
  \begin{subfigure}{0.49\textwidth}\includegraphics[width=\linewidth]{capturas/02-payday-factores.png}\caption{}\end{subfigure}\hfill
  \begin{subfigure}{0.49\textwidth}\includegraphics[width=\linewidth]{capturas/04-deadspace-factores.png}\caption{}\end{subfigure}
  \begin{subfigure}{0.49\textwidth}\includegraphics[width=\linewidth]{capturas/02-payday-motivos.png}\caption{}\end{subfigure}\hfill
  \begin{subfigure}{0.49\textwidth}\includegraphics[width=\linewidth]{capturas/04-deadspace-motivos.png}\caption{}\end{subfigure}
  \caption{El recorrido en la ficha: PAYDAY 3 a la izquierda y Dead Space a la derecha. La banda
  (a, b), lo que mueve la estimación (c, d) y los motivos más frecuentes (e, f). Capturas
  provisionales del sitio: se regeneran desde el tag de entrega.}
  \label{fig:recorrido}
\end{figure}

\begin{figure}[htbp]
  \centering
  \begin{subfigure}{\textwidth}\centering\includegraphics[width=0.86\linewidth]{capturas/05-comparar.png}\caption{}\end{subfigure}
  \begin{subfigure}{0.49\textwidth}\includegraphics[width=\linewidth]{capturas/06-explorar-accion-bajo.png}\caption{}\end{subfigure}\hfill
  \begin{subfigure}{0.49\textwidth}\includegraphics[width=\linewidth]{capturas/07-inicio-reembolso.png}\caption{}\end{subfigure}
  \caption{Tres accionables: Comparar (a), el estante de riesgo bajo de Acción en Explorar (b) y la
  ventana de reembolso en el Inicio (c), cuyos motivos se cuentan sobre todo el catálogo servido.
  Capturas provisionales del sitio: se regeneran desde el tag de entrega.}
  \label{fig:accionables}
\end{figure}

\subsection{Usabilidad}\label{sec:estrategia-usabilidad}

La banda se dice con palabras y color, no con un porcentaje: el score no es una probabilidad
calibrada (sección~\ref{sec:modelacion-umbrales}), y un porcentaje se leería como una. La ficha
explica la banda con sus tres factores principales, cada uno con su nivel de evidencia, y los
motivos con el número de reseñas sobre el que se calcularon, para que una cifra sobre pocas reseñas
no parezca más firme de lo que es. Todo lo que muestra el riesgo es igual para cualquier persona,
así que dos personas que consultan el mismo juego ven lo mismo. Los controles pequeños miden al
menos \cifraObjetivoTactil{} píxeles de alto, para tocarlos con el dedo en un teléfono, y el sitio tiene tema claro y
oscuro.
